% =============================================================================
% cap3_pipeline/06_sintese_ponte_cap4.tex
% Seção 3.6 -- Síntese do pipeline e ponte com a metodologia
%   v2.2 (quarta revisão): enxuta; escopo, métricas e testes movidos ao Cap. 4.
% =============================================================================
\section{Síntese e ponte com a metodologia}
\label{sec:pipeline-sintese}

As quatro fases descritas neste capítulo formam um pipeline cuja lógica de conjunto é mais importante que qualquer um de seus modelos individuais. A Fase~I (Seção~\ref{sec:pipeline-fase1}) absorve as mídias entregues pela tipologia operacional --- vídeo H.264/H.265 gravado em cartão microSD pela IPCAM adaptada --- e a converte em um conjunto canônico de quadros anotados com metadados de proveniência por ponto. A Fase~II (Seção~\ref{sec:pipeline-fase2}) usa o \modeloMD~como filtro generalista para descartar agressivamente \emph{ghost frames} e quadros vazios, comprimindo o volume de entrada em uma ou duas ordens de grandeza. A Fase~III (Seção~\ref{sec:pipeline-fase3}) emprega o \modeloSN~para discriminar a espécie do animal detectado, filtrando a fauna não-alvo característica do Campus~2 e encaminhando apenas registros confiantes de \estr{Felis catus} adiante. A Fase~IV (Seção~\ref{sec:pipeline-fase4}) executa a reidentificação individual com \modeloPetFace~e com o método de partes do corpo de~\citeonline{akbar2025bodypart}, combinando face e flanco em decisão ranqueada contra o catálogo da colônia, com curadoria humana acionada nos casos duvidosos. O resultado é a transformação de uma massa bruta de mídia em série temporal estruturada de \emph{visitas individuais por ponto}, que é exatamente a informação necessária para os fins de manejo declarados no Capítulo~\ref{cap:introducao}.

Em coerência com a decisão arquitetural por tipologia única (Seção~\ref{sec:decisao-arquitetural}), o pipeline é executado integralmente em servidor único, sobre janelas históricas já consolidadas pelo ciclo \emph{offline-first} de retirada de cartões. Nenhuma das quatro fases pressupõe processamento em borda, e a sequência de execução é estritamente serial. A possibilidade de deslocar a Fase~II para a borda permanece como otimização opcional, retomada como trabalho futuro no Capítulo~\ref{cap:conclusao}; sua eventual adoção alteraria somente \emph{onde} a Fase~II é executada, não \emph{como} o restante do pipeline opera. Essa simplicidade arquitetural mantém o sistema reprodutível em uma única máquina e está alinhada à escala e ao prazo de um projeto de graduação.

Com o pipeline assim consolidado em quatro fases e em uma única máquina, o Capítulo~\ref{cap:metodologia} desenvolve o protocolo de avaliação correspondente, restrito a datasets públicos já disponíveis e organizado em torno das quatro fases aqui descritas; o Capítulo~\ref{cap:resultados} apresenta os resultados preliminares dessa avaliação e o Capítulo~\ref{cap:conclusao} retoma as limitações e os trabalhos futuros decorrentes, em particular a coleta sistemática de dados no Campus~2.
